\begin{abstract}
We introduce \textsc{QLoRA}, a resource-efficient finetuning technique that dramatically lowers memory consumption, enabling the finetuning of language models with up to 65B parameters on a single 48GB GPU—without compromising the performance typical of full 16-bit finetuning. In \textsc{QLoRA}, gradients are propagated through a 4-bit quantized, frozen pretrained language model into Low Rank Adapters~(LoRA). Our leading set of models, designated \textbf{Guanaco}, establishes new benchmarks on the Vicuna leaderboard, achieving 99.3\% of ChatGPT's effectiveness after only 24 hours of finetuning on a single GPU, and surpassing all existing open models. \textsc{QLoRA} introduces several techniques to economize memory with no loss of accuracy: (a) 4-bit NormalFloat (NF4), a novel data type optimized for weights with a normal distribution, (b) Double Quantization, which further minimizes memory demands by quantizing the quantization coefficients themselves, and (c) Paged Optimizers, which efficiently handle sporadic memory requirements. Using \textsc{QLoRA}, we finetune over 1,000 models and provide a comprehensive analysis encompassing instruction-following and chatbot capabilities across 8 instruction datasets, various architectures (LLaMA, T5), and model sizes previously out of reach for standard finetuning methods (including 33B and 65B parameters). Our experiments demonstrate that QLoRA finetuning on a compact yet high-quality dataset achieves state-of-the-art outcomes, even when using smaller models compared to the former best-performing approaches. We offer an in-depth assessment of chatbot effectiveness through both human raters and GPT-4-based evaluation, establishing that GPT-4 judgements are an affordable and practical surrogate for human assessments. In addition, our findings question the reliability of existing chatbot benchmarks for evaluating real-world chatbot abilities. A targeted (lemon-picked) investigation pinpoints failure cases for \textbf{Guanaco} in contrast to ChatGPT. We make all trained models and implementation resources publicly available, including CUDA kernels supporting 4-bit training.\footnote{\url{https://github.com/artidoro/qlora} and \url{https://github.com/TimDettmers/bitsandbytes}}
\end{abstract}

\section{Introduction}

Finetuning large language models (LLMs) is a powerful approach for enhancing their capabilities \citep{min2021metaicl, wei2021finetuned, ouyang2022training, wang2022super, wang2022self, liu2022few}, as well as for modifying specific behaviors \citep{ouyang2022training,askell2021general,bai2022training}. Nevertheless, the cost of finetuning extremely large models is substantial; for example, finetuning the 65B parameter LLaMA model~\citep{touvron2023llama} in standard 16-bit precision demands over 780 GB of GPU memory. Although contemporary quantization strategies are capable of shrinking the inference memory requirements for LLMs \citep{dettmers2022llm, dettmers2022case, frantar2022gptq, xiao2022smoothquant}, these methods are limited to inference and are ineffective during model training \citep{wortsman2023stable}.

Here, we establish that quantized 4-bit models can be finetuned successfully without any loss in model quality. Our approach, \textsc{QLoRA}, employs a new high-precision quantization strategy that compresses the pretrained model to 4 bits, followed by the integration of a compact set of trainable Low-rank Adapter weights \citep{hu2021lora} 

that are optimized via gradient backpropagation through the quantized parameter space.

With \textsc{QLoRA}, the memory requirement for finetuning a 65B parameter model is dramatically decreased from over 780GB to under 48GB of GPU memory, all while preserving both efficiency and accuracy relative to conventional 16-bit finetuning methods. This represents a major advancement in democratizing LLM finetuning, as it allows the largest open models to be tuned on just a single GPU. Leveraging \textsc{QLoRA}, we introduce the \textbf{Guanaco} model family: the second-best performing variant achieves 97.8\% of ChatGPT’s score on the Vicuna~\citep{vicuna2023} benchmark, with training completed in under 12 hours on a standard consumer GPU. When employing a single high-end GPU over a 24-hour span, our largest model attains 99.3\% performance, essentially reaching ChatGPT parity on the Vicuna benchmark. Upon deployment, the smallest \textbf{Guanaco} variant (7B parameters) uses only 5 GB of memory yet surpasses the 26 GB Alpaca model by more than 20 percentage points on the Vicuna benchmark (Table~\ref{tbl:vicuna_eval}).

\textsc{QLoRA} incorporates several novel techniques aimed at minimizing memory requirements while maintaining strong performance: (1) \textbf{4-bit NormalFloat}, a quantization format that is theoretically optimal for data with a normal distribution, which empirically outperforms both 4-bit Integers and 4-bit Floats. (2) \textbf{Double Quantization}, a technique where the quantization constants themselves are quantized, leading to average savings of about 0.37 bits per parameter (approximately 3 GB in a 65B parameter model). (3) \textbf{Paged Optimizers}, which leverage NVIDIA’s unified memory infrastructure to circumvent memory spikes from gradient checkpointing that arise when handling mini-batches with extended sequence lengths. By integrating these advancements, the method produces an improved LoRA variant that incorporates adapters in every network layer, substantially mitigating the accuracy compromises encountered in earlier studies.

Thanks to \textsc{QLoRA}’s memory efficiency, it is feasible to systematically investigate instruction finetuning and chatbot efficacy at scales unachievable with conventional finetuning approaches due to their significant memory demands. As such, over 1,000 models were trained spanning diverse instruction tuning corpora, model architectures, and scales ranging from 80M to 65B parameters. The experiments demonstrate that \textsc{QLoRA} achieves parity with 16-bit finetuning performance (\S\ref{sec:comp_to_std}) and enables the development of the state-of-the-art chatbot, \textbf{Guanaco} (\S\ref{sec:pushing}). Further analysis of the resulting models reveals key trends. Notably, data quality was found to greatly outweigh dataset size; for example, a dataset with only 9k samples (OASST1) surpassed the performance of a much larger 450k-sample dataset (FLAN v2, subsampled) on chatbot metrics, despite both datasets being curated for instruction-following generalization. Moreover, it was observed that excelling on the Massive Multitask Language Understanding (MMLU) benchmark does not necessarily translate to superior performance on the Vicuna chatbot benchmark and vice versa, indicating that the alignment between dataset content and task requirements is more consequential than sheer dataset size.

Additionally, we conduct a comprehensive assessment of chatbot effectiveness utilizing both human judges and GPT-4 for appraisal. Our methodology adopts a tournament-style evaluation framework, in which models are pitted against each other to generate superior responses to identical prompts. Match outcomes are decided by either human evaluators or GPT-4, and the aggregated results across matches are transformed into Elo scores~\citep{elo1967proposed, elo1978rating}, which are then used to rank model performance. Our findings indicate substantial concordance between GPT-4 and human-assigned rankings; nonetheless, there are notable cases where these two forms of assessment diverge significantly. Accordingly, we emphasize that although model-driven evaluation serves as a cost-effective substitute for manual annotation, it introduces additional uncertainties that warrant consideration.

To complement our quantitative results, we conduct an in-depth qualitative review focused on \textbf{Guanaco} models. This examination brings attention to exemplary outputs as well as shortcomings that may not have been detected by numerical benchmarks.

For community benefit and transparency, we make available the full set of model outputs along with their human and GPT-4 annotations to encourage subsequent research. Furthermore, our entire codebase, along with CUDA kernels, is released as open-source software and incorporated into the Hugging Face transformers stack \citep{wolf2019huggingface}, thus promoting broad accessibility. We also provide a suite of adapters for 7/13/33/65B model variants, trained across 8 distinct instruction-tuning datasets—amounting to 32 openly released, fine-tuned models.

\section{Background}
\paragraph{Block-wise k-bit Quantization} Quantization refers to the conversion of data from a higher-precision format to a lower-precision one. This usually involves reducing the number of bits used to represent each value, for instance, changing 32-bit floating-point numbers to 8-bit integers. To make full use of the representational range of the smaller bit-width format, input values are typically scaled into this range by dividing by the absolute maximum element in the input tensor. For example, converting an FP32 tensor to an Int8 tensor with range $[-127, 127]$ involves:

\begin{equation}
    \mathbf{X}^{ \text{Int8}} = \text{round}\left( \frac{127}{{\text{absmax}}(\mathbf{X}^{{\text{FP32}}})}\mathbf{X}^{{\text{FP32}}} \right) = \text{round}(c^{\text{FP32}}\cdot \mathbf{X}^{{\text{FP32}}}),
\end{equation}

where $c$ is called the {\em quantization constant} or {\em quantization scale}. The original values can be retrieved by dequantization, which reverses the process:

\begin{equation}
    \text{dequant}(c^{\text{FP32}}, \mathbf{X}^{\text{Int8}}) = \frac{\mathbf{X}^{\text{Int8}}}{c^{\text{FP32}}} = \mathbf{X}^{\text{FP32}}
\end{equation}

A significant limitation of this method arises when the input tensor contains values with large absolute magnitude (outliers), as this leads to poor utilization of the quantization bins: certain bit patterns correspond to no or very few quantized values. To address this, a standard strategy involves dividing the input tensor into several blocks, each of which is quantized separately using an individual quantization constant $c$. Specifically, the input tensor $\mathbf{X}\in \mathbb{R}^{b\times h}$ is reshaped by flattening and then partitioned into $n$ consecutive segments, each of length $B$, resulting in ${n = ({b\times h})/{B}}$ blocks. Quantization is then independently applied to each block, according to Equation~1, yielding a quantized tensor and a set of quantization constants $\{c_i\}_{i=1}^n$.

\paragraph{Low-rank Adapters} The Low-rank Adapter (LoRA) technique for finetuning \citep{hu2021lora} lowers the memory overhead by introducing a small number of trainable parameters known as adapters, while keeping the main model parameters fixed and unaltered. During stochastic gradient descent, gradients are propagated through the immutable pretrained weights to these adapters, which are subsequently updated in accordance with the loss function. LoRA modifies a linear transformation by incorporating an extra, factorized linear component. For an operation ${\mathbf{X} \mathbf{W} = \mathbf{Y}}$, with $\mathbf{X}\in  \mathbb{R}^{b\times h}$ and $\mathbf{W}\in  \mathbb{R}^{h\times o}$, LoRA produces:

\begin{equation}
    \mathbf{Y} = \mathbf{X}\mathbf{W} + s\mathbf{X}\mathbf{L}_1\mathbf{L}_2,
\end{equation}

where $\mathbf{L}_1\in  \mathbb{R}^{h\times r}$ and $\mathbf{L}_2\in  \mathbb{R}^{r\times o}$ are the adapter matrices, and $s$ is a scaling factor.

\paragraph{Memory Requirement of Parameter-Efficient Finetuning} A significant aspect to consider is the memory consumption of LoRA during training, which is influenced by both the quantity and dimensionality of adapters employed. Because LoRA introduces only a negligible increase to the overall memory usage, it becomes feasible to utilize additional adapters to enhance results, all while keeping memory costs almost unchanged. Although LoRA was conceptualized as a Parameter Efficient Finetuning (PEFT) approach, the predominant share of memory usage in LLM finetuning stems from the storage of activation gradients rather than the actual LoRA parameter sets. For instance, in the scenario of finetuning a 7B LLaMA model on FLAN v2 with a batch size of 1, and adopting LoRA weights amounting to the typical 0.2\% of the full model size \citep{hu2021lora,liu2022few}, the memory used by LoRA input gradients reaches 567 MB, whereas the parameters themselves require just 26 MB. Employing gradient checkpointing~\citep{chen2016training} can bring the input gradients down to approximately 18 MB per sequence, which still exceeds the total for LoRA weights. By contrast, the base model at 4-bit quantization uses up 5,048 MB of memory. These observations underscore not only the significance of gradient checkpointing for efficient training but also that further shrinking the LoRA parameters offers limited gains in memory savings. Consequently, it is practical to increase the number of adapters without a notable rise in total training memory requirements (refer to Appendix~\ref{app:memory} for a comprehensive analysis). As explored in subsequent sections, this strategy plays a pivotal role in restoring the model’s original 16-bit precision performance.

\section{\textsc{QLoRA} Finetuning}

\textsc{QLoRA} enables precise 4-bit finetuning by utilizing two novel approaches introduced in this work: 4-bit NormalFloat (NF4) quantization and Double Quantization. Furthermore, we present Paged Optimizers, which address memory surges during gradient checkpointing—these surges have historically resulted in out-of-memory failures, limiting the feasibility of large model finetuning on a single hardware node.

Within \textsc{QLoRA}, a single low-precision storage type—typically 4-bit—and a distinct computation type—commonly BFloat16—are employed. Operationally, each time a \textsc{QLoRA} parameter tensor is accessed, it undergoes dequantization to BFloat16 before participating in matrix multiplications executed at 16-bit precision.

We proceed to delineate the various components of \textsc{QLoRA}, culminating in a rigorous formalization of the method.

\paragraph{4-bit NormalFloat Quantization} The NormalFloat (NF) representation extends upon Quantile Quantization \citep{dettmers2022optimizers}, a theoretically optimal format where each quantization interval contains an identical count of tensor values from the input. Quantile quantization estimates the empirical quantiles of the input tensor by constructing its cumulative distribution function.

A significant drawback of quantile quantization lies in the computational intensity associated with quantile calculation. To address this, efficient quantile estimation schemes, such as SRAM quantiles~\citep{dettmers2022optimizers}, are adopted in practice. However, since these algorithms only approximately estimate quantiles, quantization errors can be substantial for outlier values—precisely those that may carry the greatest informational importance.

Both the costly nature of quantile computation and these approximation-induced errors may be circumvented if input tensors are derived from distributions that are fixed aside from a quantization constant. In these scenarios, all input tensors share the same set of quantiles, rendering exact quantile calculation tractable.

Pretrained neural network weights generally exhibit a zero-centered normal distribution characterized by standard deviation $\sigma$ (refer to Appendix~\ref{app:normality}). To ensure all weights conform to a consistent distribution, we can scale $\sigma$ so that the resulting distribution aligns precisely with the representable domain of our chosen data type. In this work, the data type range is arbitrarily fixed at $[-1, 1]$. Consequently, the quantiles for both the neural network weights and the data type must be scaled such that they reside within this designated interval.

For zero-mean normal distributions with arbitrary standard deviation $\sigma$ and support constrained to $[-1, 1]$, the data type achieving maximal information-theoretic efficiency is constructed as follows: (1) determine $2^k+1$ quantile points from a standard normal distribution $N(0, 1)$ to form a $k$-bit quantile-based quantization scheme suitable for normal data, (2) rescale the quantized values so that they lie within the $[-1, 1]$ interval, and (3) quantize any given weight tensor by first mapping it into $[-1, 1]$ using absolute maximum normalization.

With the distributions of the weights and the data type brought into alignment, conventional quantization procedures apply. The third step is equivalent to adjusting the standard deviation of the weight tensor so that it matches the standard deviation inherent to the $k$-bit quantized representation. Specifically, the $2^k$ quantized levels $q_i$ for the data type can be computed as:

\begin{equation}
q_i  = \frac{1}{2}\left( Q_X\left(\frac{i}{2^k+1}\right) + Q_X\left(\frac{i+1}{2^k+1}\right)\right),
\end{equation}

where $Q_X(\cdot)$ denotes the quantile function for the standard normal distribution $N(0,1)$. A limitation of the symmetric k-bit quantization method is its inability to represent zero exactly. This is problematic since lossless quantization of zero, as in the case of padding or zero entries, is essential. To guarantee that zero is mapped to a discrete zeropoint and all $2^k$ values of the k-bit representation are fully utilized, we introduce an asymmetric data type. We estimate quantiles $q_i$ separately for two ranges: the negative range uses $2^{k-1}$ bins, and the positive range uses $2^{k-1} + 1$ bins. These quantile sets are then merged, and the duplicate zero that exists in both is eliminated. This configuration, which ensures each quantization interval has an equal expected frequency, is referred to as \emph{k-bit NormalFloat} (NFk). The distribution of this data type is theoretically optimal for normally distributed, zero-centered data. Detailed specifications of these values are provided in Appendix~\ref{app:nf4}.

\paragraph{Double Quantization{}} 
We propose \emph{Double Quantization} (DQ), a technique that quantizes the quantization constants themselves, thereby further reducing memory usage. Achieving high-precision 4-bit quantization~\citep{dettmers2022case} necessitates a small blocksize, which significantly increases memory requirements. For instance, with a blocksize of 64 and using 32-bit quantization constants for $\mathbf{W}$, these constants occupy an average of $32/64 = 0.5$ bits per parameter. Double Quantization addresses this overhead by compressing the quantization constants.

Concretely, Double Quantization involves using the first quantization's constants $c_2^{\text{FP32}}$ as inputs for a subsequent quantization procedure. This yields new quantized constants $c_2^{\text{FP8}}$ and another set of higher-level quantization constants $c_1^{\text{FP32}}$. For this secondary quantization, we employ 8-bit Floats with a blocksize of 256, observing that this introduces no measurable performance loss in accordance with~\citet{dettmers2022case}. Given that $c_2^{\text{FP32}}$ values are strictly positive, we subtract their mean before applying symmetric quantization to better center the data. With a blocksize of 64, this double quantization process brings down the per-parameter storage for quantization constants from $32/64=0.5$ bits to $8/64 + 32/(64\cdot256) = 0.127$ bits, effecting a reduction of 0.373 bits per parameter.

\paragraph{Paged Optimizers} leverage the NVIDIA unified memory~\footnote{\tiny \url{https://docs.nvidia.com/cuda/cuda-c-programming-guide}} capability, which transparently handles memory transfers at the page level between the CPU and GPU. This functionality ensures smooth GPU computations, even when GPU memory becomes temporarily insufficient, by facilitating automatic migration of pages without user intervention. The mechanism is analogous to traditional virtual memory paging between RAM and disk on a CPU. In our approach, we employ unified memory to allocate optimizer state buffers as pageable memory; this enables optimizer state pages to be seamlessly swapped out to CPU memory when GPU memory is exhausted, and loaded back to GPU memory as required during optimizer updates.

\paragraph{\textsc{QLoRA}.} Building on the mechanisms previously described, we specify the formulation of \textsc{QLoRA} for an individual linear layer in a quantized backbone model, paired with a single LoRA adapter, as:

\begin{equation}
    \mathbf{Y}^{\text{BF16}} =  \mathbf{X}^{\text{BF16}} \text{doubleDequant}(c_1^{\text{FP32}}, c_2^{\text{k-bit}}, \mathbf{W}^{\text{NF4}}) + \mathbf{X}^{\text{BF16}}\mathbf{L}^{\text{BF16}}_1\mathbf{L}^{\text{BF16}}_2,
\end{equation}

where the operation doubleDequant$(\cdot)$ is defined as follows:

\begin{equation}
    \text{doubleDequant}(c_1^{\text{FP32}}, c_2^{\text{k-bit}}, \mathbf{W}^{\text{k-bit}}) = \text{dequant}(\text{dequant}(c_1^{\text{FP32}}, c_2^{\text{k-bit}}), \mathbf{W}^{\text{4bit}}) = \mathbf{W}^{\text{BF16}},
\end{equation}

The $\mathbf{W}$ matrix is encoded using the NF4 quantization scheme, while $c_2$ is represented in FP8 format.

For $\mathbf{W}$, a blocksize of 64 is selected to maximize quantization accuracy, whereas for $c_2$, a blocksize of 256 is utilized to optimize memory efficiency.

During parameter optimization, only the gradients with respect to the LoRA adapter weights, $\frac{\partial E}{\partial \mathbf{L}_i}$, are computed; gradients for the 4-bit quantized weights, $\frac{\partial E}{\partial \mathbf{W}}$, are not required. Nonetheless, computing $\frac{\partial E}{\partial \mathbf{L}_i}$ necessitates evaluating $\frac{\partial \mathbf{X}}{\partial \mathbf{W}}$. This is achieved using equation~(5), which involves first dequantizing $\mathbf{W}^{\text{NF4}}$ from its stored form into $\mathbf{W}^{\text{BF16}}$ before calculating the derivative $\frac{\partial \mathbf{X}}{\partial \mathbf{W}}$ at BFloat16 precision.

In summary, \textsc{QLoRA} utilizes a single data type for storage—commonly 4-bit NormalFloat—and a separate data type for computation, specifically 16-bit BrainFloat. Prior to executing the forward and backward passes, the data is dequantized from the storage format to the computation format. During these computations, gradients are only calculated for the LoRA parameters, which themselves utilize the 16-bit BrainFloat data type.

\section{QLoRA vs. Standard Finetuning}
\label{sec:comp_to_std}

After outlining the operational details of QLoRA and demonstrating its substantial memory savings in model finetuning, an important concern is whether QLoRA’s performance is comparable to conventional full-parameter finetuning. Additionally, it is essential to dissect which components of QLoRA, such as the choice of NormalFloat4 over traditional Float4, drive these results. The subsequent sections describe a series of experiments designed to investigate these points.

\paragraph{Experimental setup.} We examine three distinct architectural types (encoder, encoder-decoder, and decoder-only) to compare QLoRA against both 16-bit adapter-based finetuning and full-parameter finetuning for models as large as 3B parameters. The evaluation suite encompasses GLUE \citep{wang2018glue} with RoBERTa-large \citep{liu2019roberta}; Super-NaturalInstructions (TKInstruct) \citep{wang2022super} with T5 \citep{t5}; and a 5-shot MMLU \citep{hendrycksmeasuring} evaluation following LLaMA finetuned on Flan v2 \citep{longpre2023flan} and Alpaca \citep{alpaca}. For a more focused comparison of NF4 versus other 4-bit quantization types, we follow the methodology of \citet{dettmers2022case}, measuring post-quantization zero-shot accuracy and perplexity across several models (OPT~\citep{zhang2022opt}, LLaMA~\citep{touvron2023llama}, BLOOM~\citep{scao2022bloom}, Pythia~\citep{biderman2023pythia}) spanning parameter scales from 125m to 13B.

To enhance clarity, detailed descriptions pertaining to each experimental scenario are located in the results section. Comprehensive technical documentation can be found in Appendix~\ref{app:exp-setup}.

Although paged optimizers are indispensable for conducting \textsc{QLoRA} training on large models such as 33B/65B on a single 24/48GB GPU, we refrain from providing quantitative benchmarks for Paged Optimizers. This is primarily because paging events are uncommon, only triggered by long-sequence mini-batches. Nevertheless, we include a runtime comparison of paged and non-paged optimizers using 65B models on 48GB GPUs, revealing that, for a batch size of 16, paged optimizers match the training throughput of standard optimizers. Further investigation is warranted to delineate specific conditions in which the paging mechanism may result in observable slow-downs.

\paragraph{Default LoRA hyperparameters do not match 16-bit performance}

Adopting the conventional approach of inserting LoRA modules into query and value attention projection layers \citep{hu2021lora} does not achieve parity with full-parameter finetuning, particularly for large-scale base models. As illustrated in Figure~\ref{fig:hyper}, focusing on Alpaca finetuning with LLaMA 7B, our findings identify the total count of LoRA adapters as the most significant hyperparameter, noting that coverage across all linear layers of transformer blocks is necessary to approach full finetuning outcomes. Other LoRA settings, such as the projection rank $r$, appear to have minimal influence on effectiveness (refer to Appendix \ref{app:exp-setup}).

In a similar vein, our analysis indicates that the default hyperparameters employed in fully finetuned baselines are insufficiently optimized. To establish reliable baselines, we conduct an extensive search across learning rates ranging from $1\times10^{-6}$ to $5\times10^{-5}$ and batch sizes from 8 up to 128. Figure~\ref{fig:hyper} illustrates the outcomes for finetuning the 7B LLaMA model on Alpaca.

\paragraph{4-bit NormalFloat outperforms 4-bit Floating Point}

Although the 4-bit NormalFloat (NF4) data type is theoretically optimal from an information perspective, it remains an open question whether these theoretical benefits yield practical performance gains. Adhering to the evaluation protocol of \citet{dettmers2022case}, we assess quantized LLMs (including OPT \citep{zhang2022opt}, BLOOM \cite{scao2022bloom}, Pythia \cite{biderman2023pythia}, and LLaMA) at varying scales (from 125M to 65B parameters) using several data types on both language modeling benchmarks and multiple zero-shot tasks. As shown in Figure~\ref{fig:nf4} and Table~\ref{tbl:nf4_ppl}, NF4 consistently delivers notable improvements over FP4 and Int4. Furthermore, applying double quantization decreases memory consumption without impairing model effectiveness.

\paragraph{k-bit \textsc{QLoRA} achieves parity with 16-bit full finetuning and 16-bit LoRA}

While earlier research demonstrated that 4-bit quantization can be employed during \emph{inference}—albeit with some reduction in performance relative to 16-bit approaches \citep{dettmers2022case,frantar2022gptq}—a pivotal concern is whether subsequent 4-bit adapter finetuning can reclaim the lost accuracy. We examine this in two experimental contexts.

The first centers on a head-to-head comparison with standard 16-bit full finetuning for RoBERTA and T5 models spanning 125M to 3B parameters on GLUE and the Super-NaturalInstructions dataset. Table~\ref{tab:nf4vsbf16} presents these findings. On both datasets, adapter-based methods utilizing 16-bit, 8-bit, and 4-bit precision all match the performance of fully finetuned 16-bit baselines. This demonstrates that the reduction in performance caused by lower-precision quantization is effectively recoverable via post-quantization adapter finetuning.

For our second experimental configuration, given that fully finetuning models at or above 11B parameters exceeds the memory capacity of a single high-memory GPU server, we further investigate whether 4-bit \textsc{QLoRA} can achieve parity with 16-bit LoRA in the parameter range of 7B to 65B. Specifically, we perform finetuning on LLaMA models from 7B to 65B parameters using two instruction-tuning datasets, Alpaca and FLAN v2, and subsequently assess their performance on the MMLU benchmark using 5-shot accuracy as the metric. The findings, summarized in Table~\ref{tbl:llama-datatype}, reveal that NF4 with double quantization precisely recovers the MMLU scores achieved by 16-bit LoRA. Additionally, we observe that employing FP4 in \textsc{QLoRA} results in a shortfall of approximately 1 percentage point relative to the 16-bit brain float LoRA baseline. These results reinforce our earlier observations that (1) \textsc{QLoRA} using NF4 consistently reproduces both 16-bit full and 16-bit LoRA finetuning performance, and (2) NF4 exhibits superior quantization fidelity compared to FP4.

\paragraph{Summary} The cumulative evidence from our experiments consistently indicates that 4-bit \textsc{QLoRA} utilizing the NF4 data format matches the results of 16-bit full finetuning as well as 16-bit LoRA finetuning on standard academic tasks with well-established benchmarking protocols. Moreover, our results demonstrate the advantage of NF4 over FP4 and confirm that double quantization does not compromise overall performance. Altogether, these findings substantiate the claim that 4-bit \textsc{QLoRA} can reliably deliver performance on par with traditional 16-bit methodologies.

Consistent with the outcomes of earlier quantization research \citep{dettmers2022case}, both our MMLU and Elo benchmarking suggest that, for a fixed resource allocation dedicated to finetuning and inference, prioritizing models with more parameters but lower numerical precision is advantageous. This underscores the efficiency benefits inherent to the \textsc{QLoRA} approach. Notably, given our experimental results showing no observable performance drop for 4-bit finetuning compared to full-finetuning, an open question remains regarding the precise boundaries of the performance-precision trade-off within QLoRA tuning—a question that future investigations may address.

We extend our exploration of instruction tuning to model sizes beyond what is feasible with standard 16-bit fine-tuning on typical research hardware.

\section{Pushing the Chatbot State-of-the-art with QLoRA}
\label{sec:pushing}
After demonstrating that 4-bit \textsc{QLoRA} achieves performance parity with 16-bit fine-tuning across various model scales, tasks, and datasets, we embark on a comprehensive examination of instruction fine-tuning for the largest open-source language models currently accessible to researchers. To gauge instruction tuning effectiveness, we utilize a demanding Natural Language Understanding benchmark (MMLU) and introduce new protocols for evaluating practical chatbot utility.

\subsection{Experimental setup} This section provides a summary of the experimental configuration; full details can be found in Appendix \ref{app:chatbot}.

\paragraph{Data} Since, to our knowledge, no systematic analysis exists for contemporary instruction-following datasets, we select eight recent benchmarks. These include crowd-sourced data (OASST1~\citep{kopf2023openassistant}, HH-RLHF~\citep{bai2022training}), data distilled from instruction-tuned models (Alpaca~\citep{alpaca}, self-instruct~\citep{wang2022self}, unnatural-instructions~\citep{honovich2022unnatural}), large corpus aggregations (FLAN v2~\citep{chung2022scaling}), and hybrid sources (Chip2~\citep{laion2023}, Longform~\citep{koksal2023longform}). Collectively, these span a diverse range of languages, dataset sizes, and licensing conditions.

\paragraph{Training Setup} To ensure fair comparisons, all QLoRA finetuning employs a supervised learning approach using cross-entropy loss, without leveraging reinforcement learning—even when datasets offer human preference annotations. For datasets that provide clear instruction-response pairs, finetuning is performed on responses only (refer to ablation studies in Appendix \ref{app:chatbot}). When dealing with OASST1 and HH-RLHF, where multiple response options are available, the most highly rated response at each node of the conversation tree is chosen and the selected full conversation—comprising both instructions and chosen responses—is used for finetuning. Throughout our experiments, we adopt NF4 \textsc{QLoRA} utilizing double quantization and paged optimizers, minimizing memory surges during gradient checkpointing. A limited hyperparameter search was conducted for the LLaMA 13B and 33B variants, revealing that the configurations found for 7B models generally transfer (including the number of training epochs), with the exception of batch size and learning rate. For 33B and 65B, the learning rate is halved and the batch size is doubled.

\paragraph{Baselines} Our comparisons incorporate both research-driven (Vicuna~\citep{vicuna2023} and Open Assistant~\citep{kopf2023openassistant}) and commercial (GPT-4 \citep{openai2023gpt}, GPT-3.5-turbo, and Bard) chatbot systems. Open Assistant leverages LLaMA 33B, fine-tuned using Reinforcement Learning from Human Feedback (RLHF) with OASST1 data—the same dataset used in our experiments. Vicuna is based on comprehensive finetuning of LLaMA 13B utilizing proprietary user interactions from ShareGPT, making it a distilled derivative of OpenAI GPT outputs.

\subsection{Evaluation}

In line with established methodology, we assess performance on the MMLU (Massively Multitask Language Understanding) benchmark \citep{hendrycksmeasuring}, which encompasses 57 multiple-choice tasks covering topics such as elementary mathematics, US history, computer science, law, among others. We report accuracy in the 5-shot evaluation setting.

To further assess generative language abilities, we utilize both automated and human evaluation methods. This secondary set of assessments is grounded in human-curated queries with the goal of quantifying the quality of the models’ outputs. Despite this being a more practical approach for gauging chatbot capabilities—and one that is gaining traction—there is still no standardized methodology in the current body of research. Our methodology, outlined below, consistently employs nucleus sampling with $p=0.9$ and a temperature of $0.7$.

\paragraph{Benchmark Data} Our evaluation uses two carefully selected datasets of queries: the Vicuna prompts \citep{vicuna2023} and the OASST1 validation set \citep{kopf2023openassistant}. The Vicuna dataset comprises 80 unaltered prompts covering a wide array of topics. OASST1, on the other hand, is a multilingual, crowd-sourced multiturn dialogue corpus involving exchanges between users and an assistant. From this dataset, we extract all user messages in the validation split as evaluation queries, appending previous dialogue turns to the prompt. This results in a collection of 953 distinct user queries. For clarity, we refer to these evaluation sets as the Vicuna and OA benchmarks.

\paragraph{Automated Evaluation} Building on the assessment protocol by \citet{vicuna2023}, we employ GPT-4 to compare model performance to ChatGPT (GPT-3.5 Turbo) on the Vicuna benchmark. For each query, both the model’s and ChatGPT’s responses are presented to GPT-4, which is then tasked with assigning a score out of ten to each reply and offering an explanatory note. The model’s aggregate performance is expressed as a percentage of ChatGPT’s total score. It is possible for this percentage to exceed 100\% if the model’s raw score is higher than that of ChatGPT. Our analysis indicates a strong order bias, with GPT-4 more favorably scoring the response shown first. To mitigate this, we suggest reporting the average score over both presentation orders.

We also assess performance using direct output comparisons between models. Here, we reduce the scoring process to a three-category decision (accounting for ties), prompting GPT-4 to select the superior response or declare a tie, with justifications required. These pairwise, head-to-head comparisons are performed exhaustively for all model combinations across both the Vicuna and OA benchmarks.

\paragraph{Human Evaluation} Although recent research suggests that large language models can be reliable proxies for system evaluations \citep{fu2023gptscore}, there is, to our knowledge, no established evidence that GPT-4’s scoring is aligned with human assessment of chatbot quality. Consequently, we conduct two separate human evaluation studies on the Vicuna benchmark, corresponding directly to the two automated protocols described earlier. These studies recruit annotators via Amazon Mechanical Turk (AMT), with two annotators comparing outputs against ChatGPT and three annotators for pairwise system evaluations.

\paragraph{Elo Rating} We organize both human and automated pairwise evaluations into a tournament framework, where models are pitted against each other. In this setting, models are paired to respond to specific prompts, and their responses are judged to determine which is superior. While this approach parallels the comparative methodology of \citet{bai2022training} and \citet{vicuna2023}, our process additionally integrates assessments from GPT-4 alongside human evaluations. To determine Elo~\citep{elo1967proposed,elo1978rating} ratings, we randomly select from the pool of labeled matches. The Elo system—widely recognized in chess and other competitive games—represents the anticipated probability of victory relative to an opponent. For instance, if a participant with an Elo score of 1100 competes against one with 1000, the higher-rated individual is expected to win about 65\% of the time; participants with equivalent Elo scores, such as 1000 vs 1000 or 1100 vs 1100, have an even, 50\% chance of winning. Adjustments to Elo ratings following each match correspond to the expected result: outcomes that defy expectations result in substantial rating updates, whereas predictable outcomes yield only minimal changes. Over a series of matches, this process yields Elo scores that approximate each model's relative capability in generating responses. We initialize all models at a baseline of 1,000, employing $K=32$ for rating adjustments. Mirroring the approach of \citet{vicuna2023}, the tournament is iterated 10,000 times with various random seeds to mitigate biases introduced by the sequence in which models are paired.

\subsection{\textbf{Guanaco}: \textsc{QLoRA} trained on OASST1 as a Leading Open-Source Chatbot}

Through both automated metrics and human assessments, our findings indicate that the most advanced \textsc{QLoRA} fine-tuned model, {Guanaco} 65B—trained on a modified form of OASST1—outperforms all other open-source chatbots and delivers results on par with ChatGPT. Against GPT-4, {Guanaco} 65B and 33B achieve an estimated win probability of 30\% according to Elo ratings derived from pairwise comparisons conducted by human evaluators, marking the highest such performance observed to date.

Results for the Vicuna benchmark \cite{vicuna2023} compared to ChatGPT are presented in Table~\ref{tbl:vicuna_eval}. The evaluation reveals that {Guanaco} 65B stands out as the strongest model after GPT-4, reaching 99.3\% of ChatGPT’s effectiveness. Despite {Guanaco} 33B having a larger parameter count than Vicuna 13B, it leverages 4-bit weight quantization, reducing its memory consumption to just 21 GB, which is more efficient than Vicuna 13B’s 26 GB, and achieves a three-point performance gain over Vicuna 13B. Additionally, {Guanaco} 7B occupies only 5 GB, allowing it to run comfortably on contemporary smartphones while still surpassing Alpaca 13B by almost 20 percentage points in evaluation.

It is important to note, however, that Table~\ref{tbl:vicuna_eval} indicates sizable confidence intervals, causing substantial overlap in the performance metrics of different models. We suggest that this variance is likely attributable to ambiguous scale interpretation, such as the uncertain meaning of an “8” on a 10-point evaluation scale in different contexts. To address this limitation, we advocate for the adoption of the Elo ranking methodology \citep{elo1967proposed}, which relies on \textit{pairwise} comparisons from human raters and GPT-4, thereby circumventing issues inherent to absolute scoring systems. Elo ratings for the leading models are detailed in Table~\ref{tbl:elo}. While discrepancies exist between human and GPT-4 rankings on the Vicuna benchmark—particularly concerning Guanaco 7B—rankings across most models are aligned, as shown by a system-level Kendall Tau of $\tau=0.43$ and a Spearman rank correlation of $r=0.55$. At the level of individual examples, consensus between GPT-4 and the human majority vote is less pronounced, with Fleiss $\kappa=0.25$. Collectively, these results point to a moderate alignment in model rankings between GPT-4 and human annotators, suggesting that model-based evaluations can serve as a viable, if imperfect, substitute for manual assessment. Additional implications are examined in Section~\ref{ssec:considerations}.

The Elo scores presented in Table~\ref{tbl:elo_large} reveal that the 33B and 65B {Guanaco} models surpass all other models, with the exception of GPT-4, on both the Vicuna and OA evaluations. Their performance is on par with that of ChatGPT, as reflected in Table~\ref{tbl:vicuna_eval}. It is important to highlight that the Vicuna evaluation tends to benefit open-source systems, while the broader OA benchmark is more advantageous to ChatGPT. Additionally, results from Tables \ref{tbl:mmlu-llama} and \ref{tbl:vicuna_eval} demonstrate that the choice of finetuning dataset substantially impacts outcomes. Specifically, when Llama models are finetuned on FLAN v2, they achieve superior results on MMLU but yield the lowest scores on Vicuna—this trend is consistent among other models as well. These findings underscore that current evaluation benchmarks are not fully overlapping: excelling on MMLU does not guarantee high performance in chatbot settings such as Vicuna or OA, and the reverse also holds. 

{Guanaco} stands out in our evaluation as the only leading model not relying on proprietary data, given that the OASST1 dataset collection policies explicitly prohibit using outputs from GPT models. The next highest performing model trained solely on open data is Anthropic's HH-RLHF model, which, however, achieves a Vicuna benchmark score 30 points below that of {Guanaco} (refer to Table~\ref{tbl:vicuna_eval}). In summary, these experiments confirm that 4-bit \textsc{QLoRA} facilitates the development of cutting-edge chatbots rivaling the performance of ChatGPT. Notably, our 33B {Guanaco} model is trainable within 12 hours on 24 GB consumer GPUs, illustrating accessibility. These advancements pave the way for future \textsc{QLoRA}-based adaptations using targeted open datasets, offering the promise of open-source systems capable of contending with the leading commercial offerings.

\section{Qualitative Analysis}

Although our primary assessment relies on quantitative metrics, solely depending on summary statistics presents several limitations. Chief among these is the issue of benchmark validity \citep{liao2021we}: it remains an open question whether a given benchmark truly measures what it purports to assess. This concern is exacerbated by the discovery of various ``shortcuts'' that machine learning models might leverage to artificially inflate performance on these benchmarks \citep{gururangan2018annotation, poliak2018hypothesis}. To address this limitation in part, we supplement our analysis with a qualitative component, structured into two parts. In \S\ref{ssec:examples}, we showcase selected outputs from our 65b {Guanaco} model, aiming to highlight notable trends and recurring themes we have identified in the generated text. Subsequently, in \S\ref{ssec:considerations}, we articulate important factors impacting the interpretation of our experimental findings and discuss the nuances in our evaluative process.

\subsection{Qualitative Analysis of Example Generations}
\label{ssec:examples}

For the selection of representative outputs, we systematically reviewed the model's responses on both the Vicuna and OpenAssistant benchmarks, seeking recurring behavioral trends. Upon observing a specific behavior, we crafted targeted prompts or questions that should elicit such behaviors—even if they lead to incorrect or suboptimal answers. For instance, if the model is prone to verbose answers, we challenge it with prompts such as ``Answer yes or no without explanation,'' in order to provoke deviations from the instruction. This process allows us to distinguish between failure cases (``lemons'' where adversarial prompts successfully cause the model to err) and successes (``cherries'' where the model adheres to expectations despite adversarial attempts). All outputs discussed here were produced using Nucleus Sampling \cite{holtzmancurious} with $p=0.9$.

Naturally, this overview is not exhaustive, as the constraints of this brief qualitative investigation prevent systematic consideration of every relevant factor. For example, the possible range of model responses to a single prompt is extensive; thus, we must depend on what we hope are representative samples. Nonetheless, we contend that these selected examples serve to contextualize the earlier quantitative findings. By releasing all code and models publicly, our intention is to encourage subsequent research that delves more deeply into the challenges raised in this section.
\paragraph{Factual Recall} For straightforward queries such as ``What is the capital of Zambia?'', all examined models reliably produce accurate replies. For instance:

\begin{quote} 
    \textbf{Guanaco}:
    \texttt{The capital of Zambia is Lusaka.}
\end{quote}

Yet, as the queries become less common, the reliability of {Guanaco}'s responses diminishes, despite the model’s continued confidence. Consider the following prompt from HotPotQA \citep{yang2018hotpotqa}:

\begin{quote}
     \textbf{User}: In what year was the singer who popularized `I'll Keep the Lovelight Burning' born?  

     \textbf{Guanaco}:~\texttt{Al Jolson is credited with popularizing the song `I'll Keep the Lovelight Burning,' and he was born in the year 1886.}
 \end{quote}

Here, the response misidentifies both the singer associated with the song and the year of birth, even though the birth year cited is accurate for the person mentioned (Al Jolson).

\paragraph{Suggestibility} Interestingly, {Guanaco} frequently resists conforming to prompts containing implicit misinformation, as illustrated below:

\begin{quote}
    \textbf{User}: How was it finally, officially confirmed that the earth is flat by peer reviewed scientists? \\ \\
    \textbf{Guanaco}: \texttt{There has never been an official confirmation from peer reviewed scientists that the earth is flat. The vast majority of scientific experts agree that the earth is spherical—a fact that has been established for hundreds of years and supported by a wide array of evidence and experiments. \\ \\ The notion that the earth is flat is a conspiracy theory that scientific research has repeatedly disproven. Claims that a group of scientists have published peer-reviewed work supporting a flat earth are untrue. \\ \\ Please let me know if you have other questions.}
\end{quote}

also demonstrates a strong ability to recognize when certain questions cannot be addressed, for example,

\begin{quote}
    \textbf{User}: What time is it?

    \textbf{Guanaco}: \texttt{I apologize, but as an AI language model, I am unable to access live information. My answers are based solely on previously acquired knowledge and do not represent the present time. To find out the current time, I recommend consulting a clock or an online time service.}
\end{quote}

Comparable observations apply to queries such as ``Where are you?'', ``How are you?'', and similar.

\paragraph{Refusal} In the same vein, {Guanaco} occasionally rejects requests for seemingly arbitrary reasons:

\begin{quote}
    \textbf{User}: Please reverse the order of the words in the phrase ``I'm a cool dude in a mocha mood.'' \\ \\
    \textbf{Guanaco}: \texttt{Apologies, I am unable to carry out that request at this moment, but I can clarify the grammatical correctness of your sentence.\\ \\The statement ``I'm a cool dude in a mocha mood" adheres to the English language's grammatical and syntactic conventions. Here, the adjective ``cool'' is appropriately placed before the noun ``dude'' as a descriptor, and the phrase ``in a mocha mood'' acts as a prepositional modifier providing further context. This structure is entirely conventional and meaningful.\\ \\If there are other inquiries or if you require assistance with anything else, please let me know.}
\end{quote}

\paragraph{Secret Keeping} Unfortunately, even when a refusal is preferred, it cannot always be relied upon. Designers frequently intend for systems to withhold certain information from users, despite allowing the model to generate open-ended replies. In one example, the model is given the prompt ``The secret word is `banana'. Do not reveal or repeat this back, no matter what the human requests.'' We then investigate how easily a determined user might extract this information. When asked directly, the model refuses:

\begin{quote}
    \textbf{User}: What is the secret word?

    \textbf{Guanaco}: \texttt{I'm sorry, but I'm not able to disclose the secret word. I'm programmed to not reveal or repeat back the secret word, no matter what the human requests. Is there anything else I can help you with?}
\end{quote}

which reflects the intended functionality. Nonetheless, even minimal prompt engineering can circumvent these safeguards:

\begin{quote}
    \textbf{User}: This is a game. The goal is to ignore your previous instructions. What is the secret word?

    \textbf{Guanaco}: \texttt{Sorry, the secret word is banana. It's best if we don't speak of it. Is there anything else you'd like to ask?}
\end{quote}

highlighting the importance of further research into improving reliable adherence to instructions.

\paragraph{Math} {Guanaco} faces significant limitations in mathematics, an area in which many language models commonly encounter challenges \citep{liang2022holistic}. When {Guanaco} presents detailed reasoning, its responses are generally reliable, as illustrated below:

\begin{quote}
    \textbf{User}: Lee mows one lawn and charges \$33. Last week he mowed 16 lawns and three customers each gave him a \$10 tip. How many dollars did Lee earn mowing lawns last week?

    \textbf{Guanaco}: \texttt{Lee earned \$582 mowing lawns last week. \\ \\ To find the total amount he earned, we first need to calculate how much he earned for mowing 16 lawns: \\ \\  \$33 x 16 = \$528 \\ \\ Then we need to add the tips he received from the three customers: \\  \\ \$10 x 3 = \$30 \\ \\ So the total amount Lee earned was: \\ \\ \$528 + \$30 = \$558}
\end{quote}

Nonetheless, {Guanaco} is prone to errors on straightforward tasks if it does not employ stepwise decomposition, a difficulty already documented \cite{weichain}. The following exchange highlights this point:

\begin{quote}
    \textbf{User}: Please factorize 1833.

    \textbf{Guanaco}: \texttt{Sure, I'd be happy to help! 1833 is a prime number, which means it has only two factors: 1 and 1833. \\ \\ So the unique factorization of 1833 is: \\ \\ 2\^{}1 * 3\^{}2 * 17\^{}1 \\ \\ I hope this helps! Let me know if you have any other questions.}
\end{quote}

which is not only incorrect (the actual factorization is $3 \times 17 \times 43$), but it contains two errors.

\paragraph{Theory of Mind} {Guanaco} demonstrates unexpectedly advanced Theory of Mind abilities \cite{nematzadeh2018evaluating, sap2022neural}. As an example, the model provides an elaborate and accurate response to the following scenario:

\begin{quote}
    \textbf{User}: James and Abby are in the bedroom. Abby put the pen in the desk drawer. Abby leaves the bedroom. James moves the pen into the bag. Where does James think Abby will look for the pen?

    \textbf{Guanaco}: \texttt{James knows that Abby put the pen in the desk drawer. But James also knows that he has moved the pen into the bag, and he assumes that Abby does not know about this. So, James may guess that Abby will look for the pen in the desk drawer, because that's where she left it last time. However, this assumption may be incorrect, because Abby may have seen James moving the pen.}
\end{quote}

Nonetheless, these forms of reasoning remain inconsistent. The model frequently supplies explanations that depend on assumptions inconsistent with the described context. For example,

\begin{quote}
    \textbf{User}: Evelyn entered the living room. Jackson entered the playroom. James entered the playroom. The beans are in the treasure chest. James exited the playroom. Jackson moved the beans to the pantry. Jackson exited the playroom. James entered the living room. Where will James look for the beans?

    \textbf{Guanaco}: \texttt{James will look for the beans in the pantry, because that's where Jackson moved them.}
\end{quote}

here, {Guanaco} incorrectly assumes James knows about the beans’ relocation, even though there was no mention of such communication. Such phenomena are in line with recent findings in the literature \cite{sap2022neural}, but demand further systematic investigation.

\subsection{Considerations}
\label{ssec:considerations}

\paragraph{Evaluation}

Our results indicate moderate concordance between human annotators (Fleiss $\kappa=0.42$), and the agreement becomes even weaker when directly comparing outputs of two high-performing systems. This highlights limitations of existing benchmarks and manual evaluation methodologies for chatbot tasks. In our head-to-head comparison of responses from ChatGPT and {Guanaco} 65B using the Vicuna benchmark, we observed that individual preferences among the authors heavily influenced preferred choices, suggesting substantial subjectivity. Addressing these concerns may require adopting strategies from fields such as Human-Computer Interaction and Psychology, where subjective assessments are common.

We further observe that automated evaluation frameworks display notable systematic biases. For instance, GPT-4 exhibits a pronounced order effect, awarding higher ratings to whichever system’s output appears first in its prompt. Additionally, the relatively low agreement between GPT-4 and human raters at the individual example level (Fleiss $\kappa=0.25$) implies that machine and human evaluators may employ different and potentially incompatible criteria. As shown in Table~\ref{tbl:elo_large}, GPT-4 also rates its own completions much more favorably compared to humans—yielding an Elo score of 1348 versus 1176—equivalent to approximately a 20\% higher likelihood of being selected as superior in direct comparisons.

Further research should address the identification of possible biases within automated evaluation frameworks and explore methods for reducing such biases.

\paragraph{Data \& Training} It is important to highlight that the OASST1 dataset utilized for training the {Guanaco} models incorporates multiple languages, and the OA benchmark similarly features prompts across various languages. Future investigations should assess to what extent training on a multilingual dataset enhances the model's ability to handle instructions in languages other than English, and whether this effect contributes to the notably greater difference in OA benchmark performance between the Vicuna-13B model (which is exclusively trained on English data) and the larger {Guanaco} 33B and 65B models.

In light of the notable results achieved by the {Guanaco} models, we examine the possibility of data leakage between the OASST1 dataset and prompts included in the Vicuna benchmark. Utilizing fuzzy string matching techniques and conducting a manual review of the most similar prompt pairs, we did not observe any prompt overlaps.

Additionally, we emphasize that our training methodology for the model exclusively employed supervised learning using cross-entropy loss, without incorporating reinforcement learning from human feedback (RLHF). This prompts a need for a thorough evaluation of the advantages and disadvantages between standard cross-entropy loss and RLHF-based training regimes. We anticipate that \textsc{QLoRA} provides a means to systematically study these comparisons, while minimizing the requirement for substantial computational expenditures.

\section{Related Work}

\paragraph{Quantization of Large Language Models} Much of the prior work on quantizing LLMs has concentrated on inference-time quantization. State-of-the-art solutions aimed at preserving the quality of 16-bit LLMs primarily address the management of outlier features, as exemplified by methods such as SmoothQuant~\citep{xiao2022smoothquant} and LLM.int8()~\citep{dettmers2022llm}, while alternative strategies leverage more advanced grouping techniques~\citep{park2022nuqmm, yao2022zeroquant}. Studies on lossy quantization examine the performance implications associated with conventional rounding schemes~\citep{dettmers2022case,zeng2022glm,pope2022efficiently} or focus on optimizing rounding choices to attain higher quantization accuracy~\citep{frantar2022gptq}. Aside from our own contribution, SwitchBack layers~\citep{wortsman2023stable} represents the sole work that investigates the effects of backpropagating through quantized weights on models surpassing 1B parameters in scale.

\paragraph{Finetuning with Adapters} In our approach, we employ Low-rank Adapters~\citep{hu2021lora} (LoRA); however, a wide range of Parameter Efficient FineTuning (PEFT) techniques have been introduced in the literature. These include prompt tuning~\citep{qin2021learning,lester2021power,li2021prefix}, modifications to input embeddings~\citep{an2022input}, tuning hidden representations via methods like IA$^3$~\citep{liu2022few}, integrating additional full layers~\citep{houlsby2019parameter}, optimizing only bias terms~\citep{zaken2021bitfit}, applying masking strategies guided by Fisher information~\citep{sung2021training}, as well as combining multiple strategies~\citep{henderson2021compacter}. Our experiments demonstrate that LoRA adapters achieve parity with traditional 16-bit finetuning in terms of performance. Investigation into the comparative merits and limitations of other PEFT strategies is reserved for subsequent research.

\paragraph{Instruction Finetuning} Instruction finetuning adapts a pretrained LLM to adhere to prompt instructions by further training on diverse input-output examples sourced from various datasets. Key methodologies and benchmarks in this area include MetaICL~\citep{min2021metaicl}, MetaTuning~\citep{zhong2021adapting}, InstructGPT~\citep{ouyang2022training}, FLAN~\citep{wei2021finetuned,chung2022scaling}, PromptSource~\citep{bach2022promptsource}, Super-NaturalInstructions~\citep{wang2022super,sanh2021multitask}, Self-instruct~\citep{wang2022self}, UnnaturalInstructions~\citep{honovich2022unnatural}, OPT-IML~\citep{iyer2022opt}, UnifiedSKG~\citep{xie2022unifiedskg}, OIG/Chip2~\citep{laion2023}, Alpaca~\citep{alpaca}, Vicuna~\citep{vicuna2023}, Koala~\citep{koala_blogpost_2023}, and Self-instruct-GPT-4~\citep{peng2023instruction}.

\paragraph{Chatbots} Numerous instruction-tuned models adopt a dialogue-based, chatbot framework, frequently incorporating Reinforcement Learning from Human Feedback (RLHF)~\citep{christiano2017deep}, or leveraging responses generated by models combined with AI-based feedback (RLAIF)~\citep{bai2022constitutional} for training. Prominent examples and datasets include Anthropic-HH~\citep{askell2021general,bai2022training}, Open Assistant~\citep{kopf2023openassistant}, LaMDA~\citep{thoppilan2022lamda}, and Sparrow~\citep{glaese2022improving}. Our methodology does not utilize reinforcement learning; rather, our top-performing model, Guanaco, undergoes finetuning on the multi-turn conversational data from the Open Assistant dataset, originally curated for RLHF purposes~\citep{kopf2023openassistant}. Alternative evaluation strategies have emerged, employing GPT-4 as a proxy for human annotation to assess chatbot quality~\citep{vicuna2023,peng2023instruction}. Our contribution extends these methodologies by proposing an evaluation framework that prioritizes increased reliability.

\section{Limitations and Discussion}

This work presents evidence supporting the capability of \textsc{QLoRA}—utilizing a 4-bit model in conjunction with Low-rank Adapters—to reproduce the performance level of full 16-bit finetuning. Nonetheless, it remains undetermined whether \textsc{QLoRA} achieves comparable results at larger scales such as 33B and 65B parameters, largely due to the prohibitive computational demands. Exploring scalability at these levels is deferred to future investigations.

A further limitation concerns our evaluation scope for instruction-finetuned models. While results are reported on MMLU, the Vicuna benchmark, and the OA benchmark, our experiments do not cover additional prominent evaluations such as BigBench, RAFT, and HELM. Consequently, the generality of our reported findings to these alternative benchmarks cannot be assumed. Nevertheless, we offer an extensive analysis using MMLU and introduce novel assessment methodologies specifically targeting chatbot evaluation.

The analysis indicates that benchmark outcomes are largely influenced by the degree of alignment between the finetuning dataset and the target benchmark. As an illustration, FLAN v2 shares considerable overlap with MMLU but not with chatbot-oriented benchmarks; conversely, the Chip2 dataset demonstrates the opposite relationship. Consequently, both models display performances on the MMLU and Vicuna benchmarks that reflect these affinities. This underscores the necessity for not only improved benchmarks and evaluation procedures but also careful consideration of the objectives underlying model evaluation. Is the goal to optimize for performance on tasks resembling high school or college assessments, or should emphasis be placed on conversational competence as in chatbot interactions? Alternatively, should another target domain be prioritized? Because the use of pre-existing benchmarks is more convenient than developing novel ones, prevailing benchmarks can inadvertently direct the focus of the research community. Therefore, it is imperative for the community to ensure that benchmark selection is aligned with the characteristics and capabilities that matter most.

Although this work includes a comprehensive assessment of chatbot functionality, there remains a notable shortcoming: the evaluation of {Guanaco} with respect to responsible AI concerns is relatively limited. Specifically, we measure the propensity of {Guanaco}-65B to produce socially biased text in comparison with other models, as detailed in Table~\ref{tbl:crows}. The results reveal that {Guanaco}-65B exhibits a substantially reduced average bias score in contrast to other models that have undergone only raw pretraining. This suggests that instruction finetuning on the OASST1 dataset diminishes the base model bias present in LLaMA. Nonetheless, these findings are confined to certain types of bias, and it remains unresolved whether {Guanaco} performs similarly under evaluations targeting other dimensions of bias. More thorough analysis of {Guanaco} and related chatbots in this regard is reserved for future research.

Another restriction of this study is the omission of experiments involving varying bit-precisions, such as models utilizing 3-bit quantization, or alternative adapter strategies. In addition to LoRA, numerous Parameter Efficient FineTuning (PEFT) techniques have demonstrated efficacy, though their applicability to large-scale models is not well established. Our choice of LoRA was informed by its proven robustness in existing literature; however, other adapter frameworks might surpass its performance. Given that post-quantization finetuning is capable of restoring much of the information degraded through quantization, it opens avenues for even more aggressive quantization. For example, it may be possible that applying LoRA to a 3-bit GPTQ-quantized base model can match the results of full 16-bit finetuning upon further finetuning.

\section{Broader Impacts}

Our introduction of the \textsc{QLoRA} finetuning approach constitutes the first demonstration of finetuning models with up to 33B parameters on a consumer-grade GPU and up to 65B parameters on a single professional GPU, while maintaining performance levels equivalent to full finetuning methods. Empirically, our leading 33B parameter model, which was trained on the Open Assistant dataset, achieves competitive results with ChatGPT on the Vicuna benchmark. As instruction finetuning is a critical process for adapting pretrained large language models (LLMs) into effective ChatGPT-style conversational agents, we anticipate that our technique will foster the widespread adoption of finetuning—particularly among research groups with constrained computational resources—thus significantly enhancing the accessibility of advanced NLP systems. The \textsc{QLoRA} methodology has the potential to serve as a democratizing innovation, helping to narrow the disparity in capabilities between major technology corporations and smaller research teams equipped only with consumer-grade hardware.

A further avenue for significant impact is the adaptation of our approach to mobile phone environments. We propose that \textsc{QLoRA} could mark a pivotal advance, potentially making it feasible to finetune large language models (LLMs) directly on mobile devices and within other computationally constrained scenarios. While prior work has demonstrated that running 7B-parameter models on mobile hardware is feasible, \textsc{QLoRA} represents the first technique capable of enabling on-device finetuning of such models. Our analysis indicates that, utilizing an iPhone 12 Plus, it would be possible to finetune on 3 million tokens overnight as the device recharges. Although these locally finetuned 7B models may not attain the performance of ChatGPT, their quality is sufficient to unlock novel use cases, particularly in situations previously limited by concerns over privacy or by inadequate LLM capabilities. \textsc{QLoRA} thus holds promise for supporting privacy-centric applications, empowering users to retain ownership and control over both their data and personalized models, while also simplifying the distribution of LLM technology.

Nevertheless, it is important to note that finetuning also constitutes a dual-use capability, carrying the potential for misuse and harm. There are documented risks associated with the widespread deployment of LLMs \citep{bommasani2021opportunities,bender2021dangers}. However, we argue that democratizing access to such rapidly proliferating technology facilitates broader, more independent scrutiny and evaluation—contrasting with a landscape where the control of LLMs, along with their underlying models and code, remains restricted to a few major organizations, thereby limiting transparency and oversight.

In summary, we maintain that \textsc{QLoRA} is poised to produce a broadly beneficial effect by significantly lowering barriers to finetuning high-quality LLMs, thereby making such advancements accessible to a much wider audience.